% !PRINTED 144
\weberpage{144}
%<<CONT
ist, und folglich ist
\weq{1}{j \geqq \frac{\Pi(n)}{[\Pi(s)]^r\,\Pi(r)}.}
Es ist leicht einzusehen, dass diese Zahl, wenn keiner der
Factoren $r, s$ gleich 1 ist, grösser als $n$ ist. Dies ergiebt sich,
wenn wir den Quotienten so schreiben:
\[
\frac{\Pi(n)}{[\Pi(s)]^r\,\Pi(r)} = \frac{(r+1)(r+2)\dots n}{(2.3\dots s)^r} =
\]
\[
\frac{n}{2}\cdot\left(\frac{r+1}{2}\cdot\frac{r+2}{2}\dots\frac{2r-1}{2}\right)\left(\frac{2r}{3}\cdot\frac{2r+1}{3}\dots\frac{3r-1}{3}\right)\dots
\]
\[
\left(\frac{(s-1)r}{s}\cdot\frac{(s-1)r+1}{s}\dots\frac{n-1}{s}\right).
\]
Denn hiernach ist
\weq{2}{j > \frac{n}{2}\left(\frac{r+1}{2}\right)^{r-1}\left(\frac{2r}{3}\right)^r \dots \left(\frac{(s-1)r}{s}\right)^r.}
Die Factoren $\dfrac{r+1}{2}, \dfrac{2r}{3}, \dots \dfrac{(s-1)r}{s}$ sind alle gleich oder
grösser als 1, denn es ist
\[
\frac{(h-1)r}{h} - 1 = \frac{(h-1)(r-1)-1}{h},
\]
also immer positiv, und der erste Factor
\[
\left(\frac{r+1}{2}\right)^{r-1}
\]
ist grösser als 2, wenn $r > 2$ ist. Ist aber $r = 2$, so ist das
Product der beiden Factoren:
\[
\left(\frac{r+1}{2}\right)^{r-1}\left(\frac{2r}{3}\right)^r = \frac{3}{2}\left(\frac{4}{3}\right)^2 = \frac{8}{3} > 2,
\]
und folglich ist unter allen Umständen $j > n$.

Ist zweitens die Gruppe $Q$ intransitiv, und zerfällt das
System der $n$ Ziffern in zwei Systeme von je $a$ und $b$ Ziffern, so
dass die Ziffern dieser beiden Systeme durch $Q$ nur unter einander
 vertauscht werden, und $n = a + b$ ist, so sind alle Permutationen
 von $Q$ in der Gruppe enthalten, die aus allen Permutationen
 der $a$ und der $b$ Ziffern besteht; d. h. der Grad
von $Q$ ist gleich oder kleiner als $\Pi(a)\Pi(b)$, und folglich der
Index $j$ von $Q$
\weq{3}{j \geqq \frac{\Pi(n)}{\Pi(a)\,\Pi(b)} = \frac{n}{1}\,\frac{n-1}{2}\dots\frac{b+1}{a}.}
